\documentclass[11pt,a4paper]{article}
\usepackage{curso}
\begin{document}
\portada{Lección}{3.8}{Longitud termodinámica y compromisos energía--tiempo}

\noindent\textbf{Unidad 3.}\quad\textbf{Objetivos:} tensor de fricción; longitud termodinámica y cota $\langle W_{\text{ex}}\rangle\ge\mathcal L^2/\tau$;
protocolo óptimo a velocidad termodinámica constante.
\textbf{Prerrequisitos:} 2.5 (fluctuación--disipación), 3.4--3.5 (límites de velocidad, diseño de protocolos).

\section{Trabajo en exceso en respuesta lineal}
Sea $V(x;\lambda)$ con parámetro de control $\lambda(t)$, $t\in[0,\tau]$, y la fuerza conjugada $X=-\partial_\lambda V$. Si el protocolo es lento
comparado con la relajación, el trabajo medio se separa en $\Delta F$ más un exceso cuadrático en la velocidad \cite{sivak}:
\begin{equation}
\langle W\rangle=\Delta F+\langle W_{\text{ex}}\rangle,\qquad
\langle W_{\text{ex}}\rangle\approx\int_0^\tau\dot\lambda\,\zeta(\lambda)\,\dot\lambda\,dt,\qquad
\zeta(\lambda)=\beta\int_0^\infty\langle\delta X(t)\,\delta X(0)\rangle_\lambda\,dt .
\end{equation}
\begin{intuicion}
$\zeta$ es la varianza de la fuerza por su tiempo de correlación: donde el sistema fluctúa mucho y se reacomoda despacio, mover el control
deja al sistema atrás y disipa. Es la misma idea del teorema de fluctuación--disipación (lección 2.5).
\end{intuicion}

\section{Longitud termodinámica}
\begin{definicion}[ (longitud termodinámica)]
$\mathcal L=\int\sqrt{\dot\lambda\,\zeta\,\dot\lambda}\,dt=\int\sqrt{\zeta(\lambda)}\,|d\lambda|$: no depende de la velocidad con que se recorre el camino.
\end{definicion}
\begin{resultado}
Por Cauchy--Schwarz, $\mathcal L^2=\big(\int_0^\tau\sqrt\zeta\,\dot\lambda\,dt\big)^2\le\tau\int_0^\tau\zeta\dot\lambda^2\,dt$, de donde
\begin{equation}
\langle W_{\text{ex}}\rangle\ \ge\ \frac{\mathcal L^2}{\tau},
\end{equation}
con igualdad si la velocidad termodinámica $\sqrt\zeta\,\dot\lambda$ es constante (geodésica a velocidad constante).
\end{resultado}

\section{Ejemplo: endurecer una trampa}
Trampa sobreamortiguada $V=\kappa x^2/2$, $\mu=\kB T=1$, $\lambda=\kappa$. Entonces $X=-x^2/2$, $\mathrm{Var}(X)=\sigma^4/2=1/(2\kappa^2)$ y $x^2$
se descorrelaciona con tasa $2\kappa$, así que $\zeta(\kappa)=1/(4\kappa^3)$ (en general $\kB T/(4\mu\kappa^3)$). Para $\kappa:1\to4$:
$\mathcal L=\int_1^4\tfrac12\kappa^{-3/2}d\kappa=1-\tfrac12=\tfrac12$, y el protocolo óptimo hace lineal a $\kappa^{-1/2}$:
$\kappa^*(t)=(1-t/2\tau)^{-2}$.
\begin{ejemplo}
Cálculo exacto con la EDO de la varianza $\dot v=-2\kappa v+2$ y $\langle W\rangle=\int\tfrac12\dot\kappa\,v\,dt$ \evS:
\begin{center}\begin{tabular}{lccccc}\toprule
$\tau$ & 1 & 2 & 4 & 8 & 16\\\midrule
rampa lineal & 0,243 & 0,140 & 0,076 & 0,041 & 0,021\\
óptimo $\kappa^*$ & 0,241 & 0,126 & 0,063 & 0,032 & 0,016\\
$\mathcal L^2/\tau$ & 0,250 & 0,125 & 0,063 & 0,031 & 0,016\\\bottomrule\end{tabular}\end{center}
Para $\tau\gtrsim2$ el protocolo óptimo satura la cota; la rampa lineal disipa $\approx20$--$30\,\%$ más. Para $\tau=1$ el régimen ya no es lineal
y la cota deja de ser exacta.
\end{ejemplo}

\begin{aplicacion}
Calendarios de recocido (lección 4.4): ir despacio donde $\zeta$ es grande, cerca de puntos críticos. Carga de acoplamientos de un TSU,
protocolos de borrado y reinicio (lecciones 1.5, 3.7) y calendarios de ruido de los modelos de difusión (Unidad 9).
\end{aplicacion}

\begin{error}
\begin{itemize}[nosep]
\item Usar $\mathcal L^2/\tau$ fuera del régimen lento. \quad Confundir $\mathcal L$ (espacio de parámetros) con $W_2$ (espacio de estados, 3.4).
\item Creer que el óptimo es lineal en $\lambda$: es constante en $\sqrt\zeta\,\dot\lambda$. \quad Olvidar que $\zeta$ crece cerca de puntos críticos.
\end{itemize}
\end{error}

\begin{thebibliography}{9}\small
\bibitem{sivak} D. A. Sivak y G. E. Crooks, Thermodynamic metrics and optimal paths, Phys. Rev. Lett. 108, 190602 (2012).
\bibitem[42]{whitelam} S. Whitelam, The thermodynamic freedom of a thermodynamic computer, 2026.
\bibitem[2]{pp} L. Peliti y S. Pigolotti, \emph{Stochastic Thermodynamics: An Introduction}, 2021.
\end{thebibliography}
\end{document}
